\begin{exercise}[Convexity of option prices in the strike]
For fixed \(S>0\), show that the maps
\[
  K\mapsto C_t(S,T,K),\qquad K\mapsto P_t(S,T,K)
\]
are convex on \((0,\infty)\), for every \(t\in[0,T]\).
\end{exercise}

\begin{solution}
Fix \(K_1,K_2>0\) and \(\lambda\in[0,1]\), and set
\[
  K_\lambda=\lambda K_1+(1-\lambda)K_2.
\]
For any terminal value \(x\), the map \(K\mapsto\pos{x-K}\) is convex.
Consequently,
\[
  \pos{S_T-K_\lambda}
  \le
  \lambda\pos{S_T-K_1}+(1-\lambda)\pos{S_T-K_2}.
\]
The right-hand side is the terminal payoff of a portfolio holding
\(\lambda\) calls with strike \(K_1\) and \(1-\lambda\) calls with strike
\(K_2\).  By the no-arbitrage comparison principle
\(\polyref{Cor.~1.1}\),
\[
  C_t(S,T,K_\lambda)
  \le
  \lambda C_t(S,T,K_1)+(1-\lambda)C_t(S,T,K_2).
\]
Thus \(K\mapsto C_t(S,T,K)\) is convex.

The same proof applies to puts, because \(K\mapsto\pos{K-x}\) is convex for
each fixed \(x\).  Therefore
\[
  \pos{K_\lambda-S_T}
  \le
  \lambda\pos{K_1-S_T}+(1-\lambda)\pos{K_2-S_T},
\]
and hence
\[
  P_t(S,T,K_\lambda)
  \le
  \lambda P_t(S,T,K_1)+(1-\lambda)P_t(S,T,K_2).
\]
This proves the two convexity statements of \(\polyref{Prop.~1.5}\).
\end{solution}
